\chapter{Complex Manifolds}\label{complex-manifolds}
\chapterstatus{complete}

\section{Introduction}\label{complex-manifolds:section-introduction}

A complex manifold is a manifold with holomorphic transition functions. The rigidity of holomorphic functions makes these spaces
very different from smooth manifolds: compact complex manifolds have only constant holomorphic functions, and their geometry is
recorded instead by the cohomology of holomorphic bundles and by the decomposition of differential forms into types, which is the
subject of Hodge theory. This chapter sets up complex manifolds, meromorphic forms on Riemann surfaces and the residue theorem on a compact Riemann surface (the sum of the residues of a
meromorphic $1$-form is zero), the uniformisation of Riemann surfaces with its consequences (the hyperbolic metric, Picard's theorem),
the operators $\partial$ and $\bar\partial$, analytic subvarieties and
blow-ups. References: \cite[Chapter 0]{griffiths-harris}, \cite[Chapter 2]{huybrechts-complex},
\cite[Chapter 3]{voisin-hodge-1}.

\noindent\textbf{Prerequisites.} Chapter~\ref{smooth-manifolds} (Smooth Manifolds), Chapter~\ref{several-complex-variables} (Several
Complex Variables), Chapter~\ref{multilinear-algebra} (Multilinear Algebra) and Chapter~\ref{vector-bundles} (Vector Bundles).

\noindent\textbf{On a first reading.} Section~\ref{complex-manifolds:section-uniformisation}, on the uniformisation of Riemann surfaces, is not used later and can be skipped.

\section{Complex manifolds}\label{complex-manifolds:section-manifolds}

\begin{definition}\label{complex-manifolds:definition-complex-manifold}
A \emph{complex manifold} of dimension $n$ is a topological $2n$-manifold with an atlas of charts $\varphi_i \colon U_i \to \CC^n$ whose transition maps are holomorphic; a map of complex manifolds is
\emph{holomorphic} if it is so in charts, and a \emph{biholomorphism} is a bijective holomorphic map with holomorphic inverse. The sheaf of holomorphic functions is $\mathcal{O}_X$, making $(X, \mathcal{O}_X)$ a
locally ringed space. Every complex manifold is in particular a smooth $2n$-manifold, canonically oriented by the orientation of $\CC^n$ (transition maps have Jacobian determinant $|\det_\CC|^2 > 0$).
\end{definition}

\begin{example}\label{complex-manifolds:example-complex-manifolds}
\begin{enumerate}
\item Open subsets of $\CC^n$; $\CC\PP^n$ with the charts of Example~\ref{smooth-manifolds:example-manifolds}, whose transition maps are rational and nowhere singular.
\item A \emph{complex torus} $X = \CC^n/\Lambda$ for a lattice $\Lambda \cong \ZZ^{2n}$; these are compact, and for $n = 1$ they are the elliptic curves. Two complex tori are rarely biholomorphic even when
they are diffeomorphic, which is the first appearance of moduli (Chapter~\ref{hodge-structures}).
\item Grassmannians $\mathrm{Gr}(k, n)$ of $k$-planes in $\CC^n$, with charts given by graphs of linear maps; they are compact complex manifolds of dimension $k(n-k)$.
\item A smooth analytic hypersurface $\{f = 0\} \subseteq \CC\PP^n$, for $f$ a homogeneous polynomial with $df \neq 0$ on the zero set, is a compact complex manifold of dimension $n - 1$, by the implicit
function theorem (Theorem~\ref{real-analysis:theorem-implicit-function}, valid for holomorphic maps by Remark~\ref{real-analysis:remark-constant-rank}).
\item Not every even-dimensional smooth manifold admits a complex structure: $S^4$ does not (Exercise~\ref{characteristic-classes:exercise-no-almost-complex}), and whether $S^6$ does is open.
\end{enumerate}
\end{example}

\begin{proposition}\label{complex-manifolds:proposition-no-functions}
A holomorphic function on a compact connected complex manifold is constant. Consequently a compact complex manifold has no nonconstant holomorphic map to $\CC^n$, and $\mathcal{O}_X(X) = \CC$.
\end{proposition}

\begin{proof}
$|f|$ attains a maximum at some $p$; in a chart around $p$, the maximum principle
(Proposition~\ref{several-complex-variables:proposition-identity-maximum}) makes $f$ constant near $p$, and the set where $f$ equals $f(p)$ is open and closed by the identity theorem.
\end{proof}

\begin{corollary}\label{complex-manifolds:corollary-compact-submanifolds}
Every compact complex submanifold of $\CC^n$ is a finite set of points. In particular no compact complex manifold of positive dimension, such as $\CC\PP^m$ or a complex torus, embeds holomorphically in $\CC^n$, in contrast with the
Whitney embedding theorem for smooth manifolds (Theorem~\ref{smooth-manifolds:theorem-whitney}).
\end{corollary}

\begin{proof}
The coordinate functions restrict to holomorphic functions on each connected component, which are constant by Proposition~\ref{complex-manifolds:proposition-no-functions}; so each component is a point, and a compact discrete set is
finite.
\end{proof}

\begin{theorem}\label{complex-manifolds:theorem-maps-of-tori}
Let $X = \CC^n/\Lambda$ and $Y = \CC^m/\Lambda'$ be complex tori. Every holomorphic map $f \colon X \to Y$ is induced by an affine map $\CC^n \to \CC^m$, $z \mapsto Az + b$, with $A$ a complex $m \times n$ matrix such that
$A\Lambda \subseteq \Lambda'$. In particular $f$ is a group homomorphism followed by a translation, and $X \cong Y$ as complex manifolds if and only if $n = m$ and $\Lambda' = A\Lambda$ for some $A \in \GL_n(\CC)$.
\end{theorem}

\begin{proof}
The projection $\CC^n \to X$ is a universal covering, as $\CC^n$ is simply connected and $\Lambda$ acts freely and properly (Theorem~\ref{topological-manifolds:theorem-covering-action}). By the lifting criterion
(Proposition~\ref{homotopy:proposition-lifting-criterion}), $f \circ \pi$ lifts to a continuous map $F \colon \CC^n \to \CC^m$, holomorphic because $\pi'$ is a local biholomorphism. For $\lambda \in \Lambda$, $F(z + \lambda) - F(z) \in \Lambda'$ for all
$z$, and by continuity and discreteness of $\Lambda'$ it is a constant $c(\lambda)$. Hence each partial derivative $\partial F_j/\partial z_k$ is $\Lambda$-periodic, so it is a holomorphic function on the compact manifold $X$, hence constant
(Proposition~\ref{complex-manifolds:proposition-no-functions}). So $F(z) = Az + b$, and $A\lambda = c(\lambda) \in \Lambda'$. If $f$ is a biholomorphism, apply this to $f$ and $f^{-1}$: their lifts compose to a lift of the identity, which is
a translation by an element of $\Lambda$, so $A$ is invertible with $A\Lambda = \Lambda'$ and $n = m$; the converse is clear.
\end{proof}

\begin{corollary}\label{complex-manifolds:corollary-elliptic-curves}
Every one-dimensional complex torus is biholomorphic to $E_\tau = \CC/(\ZZ + \ZZ\tau)$ with $\operatorname{Im}\tau > 0$, and $E_\tau \cong E_{\tau'}$ if and only if $\tau' = \frac{a\tau + b}{c\tau + d}$ for some
$\begin{pmatrix}a & b\\ c & d\end{pmatrix} \in \SL_2(\ZZ)$.
\end{corollary}

\begin{proof}
A lattice $\ZZ\omega_1 + \ZZ\omega_2$ is carried by $z \mapsto z/\omega_1$ to $\ZZ + \ZZ\tau$ with $\tau = \omega_2/\omega_1 \notin \RR$, and replacing $\omega_2$ by $-\omega_2$ if necessary, $\operatorname{Im}\tau > 0$. By
Theorem~\ref{complex-manifolds:theorem-maps-of-tori}, $E_\tau \cong E_{\tau'}$ if and only if $\alpha(\ZZ + \ZZ\tau) = \ZZ + \ZZ\tau'$ for some $\alpha \in \CC^\times$. Then $\alpha\tau = a\tau' + b$ and $\alpha = c\tau' + d$ for an integer matrix,
invertible over $\ZZ$ because the inverse change of basis is also integral; so $\tau = \frac{a\tau' + b}{c\tau' + d}$, and the determinant is $+1$ because both $\tau$ and $\tau'$ lie in the upper half plane, as
$\operatorname{Im}\frac{a\tau' + b}{c\tau' + d} = \frac{(ad - bc)\operatorname{Im}\tau'}{|c\tau' + d|^2}$. Conversely such a matrix gives the lattice isomorphism. The relation is symmetric, so $\tau'$ is also of this form in terms of $\tau$.
\end{proof}

\begin{counterexample}\label{complex-manifolds:counterexample-tori}
Diffeomorphic complex manifolds need not be biholomorphic: all $E_\tau$ are diffeomorphic to $S^1 \times S^1$, but by Corollary~\ref{complex-manifolds:corollary-elliptic-curves} $E_i \not\cong E_{2i}$, since
$\frac{ai + b}{ci + d} = 2i$ with integers of determinant $1$ would give, comparing imaginary parts, $|ci + d|^2 = \frac12$, impossible for integers $c$, $d$.
\end{counterexample}

\begin{proposition}\label{complex-manifolds:proposition-normal-form}
Let $f \colon X \to Y$ be a nonconstant holomorphic map of connected Riemann surfaces (complex manifolds of dimension $1$), and $p \in X$. There are holomorphic coordinates $z$ centred at $p$ and $w$ centred at $f(p)$ in
which $f$ is $z \mapsto z^k$ for a unique integer $k = e_p \geq 1$, the \emph{ramification index}. Consequently $f$ is open, its fibres are discrete, and if $X$ is compact and $Y$ connected, then $f$ is surjective and
has finite fibres.
\end{proposition}

\begin{proof}
Choose coordinates $z$ at $p$ and $w$ at $f(p)$ with $z(p) = 0$, $w(f(p)) = 0$. In them $f$ is a holomorphic function of one variable with $f(0) = 0$, not identically zero near $0$ by the identity theorem
(Theorem~\ref{complex-analysis:theorem-identity}) and connectedness of $X$. So $f(z) = z^kg(z)$ with $k \geq 1$ and $g(0) \neq 0$. On a small disc where $g$ has no zeros, $g'/g$ has a primitive $L$
(Corollary~\ref{complex-analysis:corollary-primitive-convex}), and $(ge^{-L})' = 0$, so $g = ce^L$ with $c \neq 0$; writing $c = e^{c'}$, the function $h = e^{(L + c')/k}$ satisfies $h^k = g$. The new coordinate
$\tilde z = zh(z)$ has derivative $h(0) \neq 0$ at $0$, so it is a local biholomorphism (Theorem~\ref{complex-analysis:theorem-open-mapping}, or
Proposition~\ref{several-complex-variables:proposition-inverse-function}), and $f = \tilde z^k$. The integer $k$ is determined by $f$, as the number of preimages near $p$ of points $w \neq 0$ near $0$. The map
$z \mapsto z^k$ is open with discrete fibres, so $f$ is. If $X$ is compact, $f(X)$ is compact, hence closed, and open, so $f(X) = Y$ when $Y$ is connected; the fibres are discrete and closed in a compact space, hence
finite.
\end{proof}

\begin{proposition}\label{complex-manifolds:proposition-degree}
Let $f \colon X \to Y$ be a nonconstant holomorphic map of compact connected Riemann surfaces. The number $\deg f = \sum_{p \in f^{-1}(q)}e_p$ does not depend on $q \in Y$. It equals the number of preimages of every
point $q$ outside the finite set $f(\{p : e_p > 1\})$ of \emph{branch points}.
\end{proposition}

\begin{proof}
The ramification points, where $f'$ vanishes in coordinates, are the zeros of a holomorphic function locally, hence discrete, hence finitely many in the compact $X$. Fix $q_0$ with preimages $p_1, \ldots, p_r$
(Proposition~\ref{complex-manifolds:proposition-normal-form}), and choose disjoint coordinate discs $D_i$ around $p_i$ on which $f$ is $z \mapsto z^{e_i}$ onto a disc $D$ around $q_0$. The set
$f(X \setminus \bigcup_iD_i)$ is compact and does not contain $q_0$, so there is a neighbourhood $D' \subseteq D$ of $q_0$ avoiding it. For $q \in D' \setminus \{q_0\}$, all preimages of $q$ lie in $\bigcup_iD_i$,
and in $D_i$ the equation $z^{e_i} = w(q) \neq 0$ has exactly $e_i$ solutions, each with ramification index $1$. So the weighted count equals $\sum_ie_i$ on $D'$. Hence $q \mapsto \deg_q f$ is locally constant, hence
constant on the connected $Y$.
\end{proof}

\begin{corollary}\label{complex-manifolds:corollary-maps-to-tori}
Every holomorphic map from a simply connected compact connected complex manifold, such as $\CC\PP^1$, to a complex torus is constant. In particular $\CC\PP^1$ is not biholomorphic, and admits no nonconstant holomorphic
map, to an elliptic curve.
\end{corollary}

\begin{proof}
Let $f \colon X \to T = \CC^m/\Lambda$. The projection $\CC^m \to T$ is a covering (Theorem~\ref{topological-manifolds:theorem-covering-action}), so by the lifting criterion
(Proposition~\ref{homotopy:proposition-lifting-criterion}) $f$ lifts to a continuous $F \colon X \to \CC^m$, which is holomorphic because the projection is a local biholomorphism. By
Proposition~\ref{complex-manifolds:proposition-no-functions} each component of $F$ is constant, so $f$ is constant. The projective line is covered by the two charts $U_0$, $U_1$, each homeomorphic to $\CC$, with
path-connected intersection $\CC^\times$; by the Seifert--van Kampen theorem (Theorem~\ref{homotopy:theorem-van-kampen}) $\pi_1(\CC\PP^1)$ is a quotient of the free product of two trivial groups, hence trivial.
\end{proof}

\section{Meromorphic forms and the residue theorem}\label{complex-manifolds:section-residues}

On a compact Riemann surface, Stokes' theorem turns the residue theorem of complex analysis into a global constraint: the residues of a meromorphic $1$-form add up to zero. This is
the simplest case of the \emph{global residue theorem} \cite[Chapter 5]{griffiths-harris}, and it has quick consequences: a meromorphic function has as many zeros as poles, a complex torus
carries no meromorphic function with a single simple pole, and the roots of a polynomial satisfy identities that are not obvious from the polynomial.

\begin{definition}\label{complex-manifolds:definition-meromorphic-form}
Let $X$ be a Riemann surface. A \emph{meromorphic $1$-form} $\omega$ on $X$ is given in every holomorphic chart $z$ by an expression $f(z)\,dz$ with $f$ meromorphic
(Definition~\ref{complex-analysis:definition-singularity}), such that for two charts $z$ and $w$ the expressions $f(z)\,dz$ and $g(w)\,dw$ agree: $g(w) = f(z(w))\,z^{\prime}(w)$. It is
\emph{holomorphic} if all the $f$ are. The \emph{order} $\operatorname{ord}_p\omega$ is $\operatorname{ord}_0f$ in a chart centred at $p$, and the \emph{residue} $\operatorname{Res}_p\omega$ is the
residue at $0$ of $f$, that is, the coefficient of $z^{-1}$ in its Laurent expansion.
\end{definition}

Since $z^{\prime}(w)$ does not vanish, the order does not depend on the chart; the residue does not either, by the lemma below. For a meromorphic function $h$, $dh = h^{\prime}(z)\,dz$ is a
meromorphic $1$-form, and so is $dh/h$ if $h \not\equiv 0$. On $\CC\PP^1$, with the charts $z$ and $w = 1/z$, the form $dz$ becomes $dz = -dw/w^2$: it has a pole of order $2$ at $\infty$ and residue $0$
there, while $dz/z = -dw/w$ has residue $1$ at $0$ and $-1$ at $\infty$.

Away from its poles, a meromorphic $1$-form is a smooth complex-valued $1$-form, $\omega = \alpha + i\beta$ with real $1$-forms $\alpha$ and $\beta$, and line integrals and integrals over manifolds
of such forms are taken componentwise. It is \emph{closed}: in a chart, with $f$ holomorphic, the Cauchy--Riemann equations (Proposition~\ref{complex-analysis:proposition-cauchy-riemann}) give
$\partial_xf = f^{\prime}$ and $\partial_yf = if^{\prime}$, so $df = f^{\prime}\,(dx + i\,dy) = f^{\prime}\,dz$ and $d(f\,dz) = df \wedge dz = f^{\prime}\,dz \wedge dz = 0$. Holomorphic charts preserve the orientation (the
Jacobian determinant of a holomorphic map is $|f^{\prime}|^2 \geq 0$), so a Riemann surface is canonically oriented, and the boundary of a coordinate disc is traversed counterclockwise in the chart.

\begin{lemma}\label{complex-manifolds:lemma-residue-integral}
Let $\omega$ be a meromorphic $1$-form on a Riemann surface $X$ and $p \in X$. For every closed coordinate disc $D$ around $p$ (the preimage of a closed disc under a holomorphic chart) whose only
possible pole is $p$,
\[
\operatorname{Res}_p\omega = \frac{1}{2\pi i}\int_{\partial D}\omega .
\]
In particular $\operatorname{Res}_p\omega$ does not depend on the chart.
\end{lemma}

\begin{proof}
If $D = \{|z| \leq r\}$ in a chart $z$ centred at $p$ with $\omega = f\,dz$, the line integral of $\omega$ over $\partial D$ is the complex path integral $\int_{|z| = r}f(z)\,dz$, which is
$2\pi i\operatorname{Res}_0f$ by the residue theorem (Theorem~\ref{complex-analysis:theorem-residues}). Now let $D^{\prime}$ be a closed coordinate disc around $p$ for another chart, with no pole other than
$p$. For $r$ small the disc $D_r = \{|z| \leq r\}$ lies in the interior of $D^{\prime}$, and $D^{\prime}$ minus the interior of $D_r$ is a compact manifold with boundary $\partial D^{\prime} \cup \partial D_r$, the
second circle with the opposite orientation. As $\omega$ is closed there, Stokes' theorem (Theorem~\ref{differential-forms:theorem-stokes}) gives $\int_{\partial D^{\prime}}\omega = \int_{\partial D_r}\omega$.
\end{proof}

\begin{theorem}[Residue theorem on a compact Riemann surface]\label{complex-manifolds:theorem-global-residue}
Let $X$ be a compact Riemann surface and $\omega$ a meromorphic $1$-form on $X$. Then $\omega$ has finitely many poles, and
\[
\sum_{p \in X}\operatorname{Res}_p\omega = 0 .
\]
\end{theorem}

\begin{proof}
In a chart the poles of a meromorphic function are isolated, so the poles of $\omega$ form a discrete closed subset of the compact $X$, hence a finite set $\{p_1, \ldots, p_k\}$. Choose pairwise
disjoint closed coordinate discs $D_j$ around the $p_j$. Then $X^{\prime} = X \setminus \bigcup_j\mathrm{int}\,D_j$ is a compact manifold with boundary $\bigcup_j\partial D_j$, each circle with the
orientation opposite to its orientation as the boundary of $D_j$, and $\omega$ is a closed form on $X^{\prime}$. By Stokes' theorem and Lemma~\ref{complex-manifolds:lemma-residue-integral},
\[
0 = \int_{X^{\prime}}d\omega = \int_{\partial X^{\prime}}\omega = -\sum_{j=1}^k\int_{\partial D_j}\omega = -2\pi i\sum_{j=1}^k\operatorname{Res}_{p_j}\omega,
\]
which is the claim.
\end{proof}

The proof is the same as that of the residue theorem in the plane, with one difference: on a compact surface there is no ``outside'', so the integral over the boundary of the complement of the
discs is not the integral over a large circle but zero. For $X = \CC\PP^1$ the point $\infty$ plays the role of the outside: the theorem says that the sum of the residues of a rational form in $\CC$
equals minus its residue at $\infty$.

\begin{corollary}\label{complex-manifolds:corollary-zeros-poles}
Let $f$ be a nonconstant meromorphic function on a compact connected Riemann surface $X$. Then $\sum_p\operatorname{ord}_pf = 0$: $f$ has as many zeros as poles, counted with multiplicity. Hence $f$
takes every value $c \in \CC$ the same number of times, counted with multiplicity, namely the number of its poles.
\end{corollary}

\begin{proof}
At a point $p$, in a chart centred at $p$, write $f = z^kg$ with $k = \operatorname{ord}_pf$ and $g$ holomorphic and nonvanishing near $0$. Then $df/f = k\,dz/z + dg/g$, and $dg/g$ is holomorphic near $0$,
so $\operatorname{Res}_p(df/f) = \operatorname{ord}_pf$. Since $f$ is not identically zero (it is nonconstant on the connected $X$, so not zero on any open set by the identity theorem), $df/f$ is a meromorphic
$1$-form, and Theorem~\ref{complex-manifolds:theorem-global-residue} gives $\sum_p\operatorname{ord}_pf = 0$. Apply this to $f - c$, which has the same poles as $f$.
\end{proof}

This is also Proposition~\ref{complex-manifolds:proposition-degree} applied to $f$ regarded as a holomorphic map to $\CC\PP^1$; the residue theorem gives it without the local analysis of that proof.

\begin{corollary}\label{complex-manifolds:corollary-no-single-pole}
Let $X$ be a compact Riemann surface with a holomorphic $1$-form $\eta$ that does not vanish at $p \in X$. Then no meromorphic function on $X$ is holomorphic on $X \setminus \{p\}$ with a simple pole at $p$.
In particular a meromorphic function on a complex torus $\CC/\Lambda$ that is not constant has at least two poles, counted with multiplicity.
\end{corollary}

\begin{proof}
Suppose $f$ were such a function. In a chart centred at $p$, $f = c/z + (\text{holomorphic})$ with $c \neq 0$ and $\eta = h(z)\,dz$ with $h(0) \neq 0$. Then $f\eta$ is holomorphic except at $p$, where
$\operatorname{Res}_p(f\eta) = c\,h(0) \neq 0$, contradicting Theorem~\ref{complex-manifolds:theorem-global-residue}. On $\CC/\Lambda$ the form $dz$ is invariant under the translations by $\Lambda$, so it
descends to a holomorphic $1$-form without zeros. A nonconstant meromorphic function has a pole, since otherwise it would be a holomorphic function on a compact connected manifold, hence constant
(Proposition~\ref{complex-manifolds:proposition-no-functions}); and a single simple pole is excluded by the first part, applied at that pole.
\end{proof}

On $\CC\PP^1$ the function $z$ has a single simple pole, at $\infty$, so by the corollary $\CC\PP^1$ has no holomorphic $1$-form that is nonzero at $\infty$. In fact it has no nonzero holomorphic
$1$-form at all: such a form is $f(z)\,dz$ with $f$ entire, and $f(z)\,dz = -f(1/w)w^{-2}\,dw$ is holomorphic at $w = 0$ only if $f(z) = O(|z|^{-2})$ as $z \to \infty$; then $f$ is bounded, hence
constant by Liouville's theorem, hence $0$.

\begin{corollary}[Euler--Jacobi formula in one variable]\label{complex-manifolds:corollary-euler-jacobi}
Let $f \in \CC[z]$ have degree $d \geq 2$ and distinct roots $r_1, \ldots, r_d$, and let $g \in \CC[z]$ have degree at most $d - 2$. Then
\[
\sum_{i=1}^d\frac{g(r_i)}{f^{\prime}(r_i)} = 0 .
\]
\end{corollary}

\begin{proof}
Apply Theorem~\ref{complex-manifolds:theorem-global-residue} to $\omega = g(z)\,dz/f(z)$ on $\CC\PP^1$. In $\CC$ its poles are the roots of $f$, which are simple since $f^{\prime}(r_i) \neq 0$, and there
$\operatorname{Res}_{r_i}\omega = \lim_{z \to r_i}(z - r_i)g(z)/f(z) = g(r_i)/f^{\prime}(r_i)$. At $\infty$, with $w = 1/z$, write $g(1/w) = w^{-e}G(w)$ and $f(1/w) = w^{-d}F(w)$ with polynomials $F$, $G$,
where $e = \deg g$ and $F(0) \neq 0$ is the leading coefficient of $f$. Then $\omega = -w^{d-e-2}\frac{G(w)}{F(w)}\,dw$, which is holomorphic at $w = 0$ because $d - e - 2 \geq 0$ (if $g = 0$ there is nothing to prove).
So the residues at the roots add up to zero.
\end{proof}

For example, for $f = z^3 - 1$ the roots are the cube roots of unity $\zeta$, with $f^{\prime}(\zeta) = 3\zeta^2$, and the formula with $g = 1$ and $g = z$ says $\sum_\zeta\zeta^{-2} = 0$ and $\sum_\zeta\zeta^{-1} = 0$;
for $g = z^2$, which has degree $d - 1$, the residue at $\infty$ is $-1$ and indeed $\sum_\zeta\frac{\zeta^2}{3\zeta^2} = 1$. For $f = z^2 - 1$ and $g = 1$: $\frac{1}{2} + \frac{1}{-2} = 0$.

In higher dimension the residue of a meromorphic $n$-form is attached to a point where $n$ divisors meet: if near $P$ the divisors are $\{f_1 = 0\}, \ldots, \{f_n = 0\}$, meeting only at $P$, and
$\omega = g\,dz_1 \wedge \cdots \wedge dz_n/(f_1 \cdots f_n)$ with $g$ holomorphic, the residue is
$\operatorname{Res}_P\omega = \frac{1}{(2\pi i)^n}\int_\Gamma\omega$ over the real $n$-cycle $\Gamma = \{|f_1| = \cdots = |f_n| = \varepsilon\}$, suitably oriented.

\begin{theorem}[Global residue theorem]\label{complex-manifolds:theorem-global-residue-higher}
Let $M$ be a compact complex manifold of dimension $n$, $D_1, \ldots, D_n$ effective divisors whose intersection is a finite set of points, and $\omega$ a meromorphic $n$-form on $M$ whose poles are
contained in $D_1 \cup \cdots \cup D_n$, of order at most one along each $D_i$ (that is, $\omega$ is a holomorphic section of $\Omega^n_M(D_1 + \cdots + D_n)$). Then
$\sum_{P \in D_1 \cap \cdots \cap D_n}\operatorname{Res}_P\omega = 0$.
\end{theorem}

This is \cite[Chapter 5]{griffiths-harris}. For $n = 1$ and $D_1$ the polar divisor it is Theorem~\ref{complex-manifolds:theorem-global-residue} for forms with simple poles. As there, the proof rests on
Stokes' theorem, now applied on the complement of tubes around the divisors, together with a local analysis of the residue. Griffiths and Harris use it to prove the Cayley--Bacharach theorem: if two plane
cubic curves meet in nine distinct points, every cubic through eight of them passes through the ninth.

\section{Uniformisation of Riemann surfaces}\label{complex-manifolds:section-uniformisation}

A \emph{Riemann surface} is a connected complex manifold of dimension $1$. The uniformisation theorem reduces the study of all Riemann surfaces to three simply connected
ones, the sphere, the plane and the disc, and to groups of automorphisms acting on them. Combined with the hyperbolic metric of the disc it shows that most Riemann surfaces,
among them all compact ones of genus at least $2$, carry a natural metric of constant curvature $-1$; and it gives a short proof of Picard's theorem. Throughout,
$\mathbb{D} = \{|z| < 1\}$ and $\mathrm{Aut}(X)$ is the group of biholomorphisms of $X$ to itself.

\begin{lemma}\label{complex-manifolds:lemma-coverings}
Let $X$ be a Riemann surface.
\begin{enumerate}
\item If $p \colon Y \to X$ is a covering map with $Y$ connected, then $Y$ has a unique structure of Riemann surface for which $p$ is a local biholomorphism, and the deck
transformations are biholomorphisms.
\item If a group $\Gamma$ acts on $X$ freely and properly by biholomorphisms, then $X/\Gamma$ has a unique structure of Riemann surface for which $X \to X/\Gamma$ is a local
biholomorphism, and a map $f \colon X/\Gamma \to Z$ to a Riemann surface is holomorphic if and only if its composite with the projection is.
\end{enumerate}
\end{lemma}

\begin{proof}
These are the proofs of Theorem~\ref{smooth-manifolds:theorem-smooth-quotient} and Example~\ref{smooth-manifolds:example-smooth-quotients}, with ``holomorphic'' in place of
``smooth'': the charts of $Y$ (of $X/\Gamma$) are the charts of $X$ composed with $p$ on sheets over evenly covered chart domains (with local inverses of the projection), and their
transition maps are transition maps of $X$, composed with the biholomorphisms of $\Gamma$ in the case of $X/\Gamma$. Topologically, $Y$ is a surface by
Proposition~\ref{topological-manifolds:proposition-constructions}(4), and $X/\Gamma$ one by (3). A deck transformation $\tau$ satisfies $p \circ \tau = p$, so in these charts it is
the identity map of coordinates.
\end{proof}

In particular the universal covering $\widetilde{X}$ of a Riemann surface (Theorem~\ref{homotopy:theorem-universal-cover}) is a simply connected Riemann surface, $\pi_1(X)$ acts
on it freely and properly by biholomorphisms as the deck group (Proposition~\ref{homotopy:proposition-deck}), and $X \cong \widetilde{X}/\pi_1(X)$.

\begin{theorem}[Uniformisation theorem]\label{complex-manifolds:theorem-uniformisation}
Every simply connected Riemann surface is biholomorphic to exactly one of $\CC\PP^1$, $\CC$ and $\mathbb{D}$.
\end{theorem}

This is due to Koebe and Poincar\'e (1907); see \cite{forster-riemann-surfaces}. The three surfaces are pairwise non-biholomorphic: $\CC\PP^1$ is compact and the others are not,
and a biholomorphism $\CC \to \mathbb{D}$ would be a bounded nonconstant entire function (Exercise~\ref{complex-analysis:exercise-disc-not-plane}). For open subsets of $\CC$ the
theorem is the Riemann mapping theorem (Theorem~\ref{complex-analysis:theorem-riemann-mapping}); the general case constructs the biholomorphism from Green's functions, using
Perron's method to solve the Dirichlet problem on the surface.

\begin{lemma}\label{complex-manifolds:lemma-automorphisms}
\begin{enumerate}
\item $\mathrm{Aut}(\CC\PP^1)$ consists of the M\"obius maps $z \mapsto (az + b)/(cz + d)$, $ad - bc \neq 0$, and each of them has a fixed point.
\item $\mathrm{Aut}(\CC)$ consists of the maps $z \mapsto az + b$, $a \neq 0$; those without fixed points are the translations $z \mapsto z + b$, $b \neq 0$.
\item $\mathrm{Aut}(\mathbb{D})$ consists of the maps $z \mapsto \lambda(z - a)/(1 - \bar az)$ with $|\lambda| = 1$ and $a \in \mathbb{D}$, and acts transitively on $\mathbb{D}$.
\end{enumerate}
\end{lemma}

\begin{proof}
(2) is Exercise~\ref{complex-analysis:exercise-automorphisms-plane}, and $az + b = z$ has the solution $b/(1 - a)$ when $a \neq 1$. (3) is
Proposition~\ref{complex-analysis:proposition-disc-automorphisms}; the map $\phi_{-a}$ sends $0$ to $a$. (1) A M\"obius map is holomorphic in the charts of $\CC\PP^1$, with the
M\"obius inverse $z \mapsto (dz - b)/(-cz + a)$, so it is an automorphism, and the M\"obius maps act transitively ($z \mapsto z + b$ and $z \mapsto 1/z$ move any point to any
other). Given $f \in \mathrm{Aut}(\CC\PP^1)$, choose a M\"obius $m$ with $m(f(\infty)) = \infty$; then $m \circ f$ restricts to an automorphism of $\CC = \CC\PP^1 \setminus \{\infty\}$,
hence is $z \mapsto az + b$ by (2), and $f = m^{-1} \circ (m \circ f)$ is M\"obius. If $c = 0$, $\infty$ is a fixed point; if $c \neq 0$, the equation $cz^2 + (d - a)z - b = 0$ has a
root $z_0 \in \CC$ (Corollary~\ref{complex-analysis:corollary-liouville}(2)), and $cz_0 + d \neq 0$, since at $z_0 = -d/c$ the left side is $(ad - bc)/c \neq 0$; so $z_0$ is a fixed point.
\end{proof}

\begin{proposition}[The Poincar\'e metric]\label{complex-manifolds:proposition-poincare-metric}
Let $g_{\mathbb{D}} = 4|dz|^2/(1 - |z|^2)^2$ on $\mathbb{D}$, that is $4(dx^2 + dy^2)/(1 - x^2 - y^2)^2$.
\begin{enumerate}
\item (Schwarz--Pick) For every holomorphic $f \colon \mathbb{D} \to \mathbb{D}$,
\[
\frac{|f^{\prime}(z)|}{1 - |f(z)|^2} \leq \frac1{1 - |z|^2},
\]
that is, $f^*g_{\mathbb{D}} \leq g_{\mathbb{D}}$: holomorphic self-maps of the disc do not increase lengths of curves, hence distances.
\item Every automorphism of $\mathbb{D}$ is an isometry of $g_{\mathbb{D}}$.
\item $g_{\mathbb{D}}$ is complete, with constant curvature $-1$.
\end{enumerate}
\end{proposition}

\begin{proof}
(1) Fix $a \in \mathbb{D}$, let $b = f(a)$ and $\phi_a(z) = (z - a)/(1 - \bar az)$. A computation gives $\phi_a^{\prime}(z) = (1 - |a|^2)/(1 - \bar az)^2$, so
$\phi_{-a}^{\prime}(0) = 1 - |a|^2$ and $\phi_b^{\prime}(b) = 1/(1 - |b|^2)$. The map $h = \phi_b \circ f \circ \phi_{-a}$ sends $\mathbb{D}$ to $\mathbb{D}$ and $0$ to $0$, so
$|h^{\prime}(0)| \leq 1$ by Schwarz's lemma (Lemma~\ref{complex-analysis:lemma-schwarz}), and by the chain rule $h^{\prime}(0) = f^{\prime}(a)(1 - |a|^2)/(1 - |b|^2)$. At a point
$z$, $f^*g_{\mathbb{D}}$ is $4|f^{\prime}(z)|^2|dz|^2/(1 - |f(z)|^2)^2$, so the inequality says $f^*g_{\mathbb{D}} \leq g_{\mathbb{D}}$, and lengths of curves decrease.

(2) Apply (1) to an automorphism and to its inverse: $f^*g_{\mathbb{D}} \leq g_{\mathbb{D}}$ and $g_{\mathbb{D}} = (f^{-1})^*f^*g_{\mathbb{D}}$ force equality.

(3) The curvature is $-1$ by Proposition~\ref{riemannian-geometry:proposition-conformal-curvature} (Example~\ref{riemannian-geometry:example-conformal-metrics}(2)). For
completeness, choose $\varepsilon > 0$ such that $\exp_0$ is defined on the ball of radius $\varepsilon$ in $T_0\mathbb{D}$. Since $\mathrm{Aut}(\mathbb{D})$ acts transitively
by isometries, and isometries map geodesics to geodesics, the same $\varepsilon$ works at every point: every geodesic of unit speed is defined on an interval of length
$\varepsilon$ around each of its times. So a maximal unit speed geodesic defined on $[0, T)$ with $T < \infty$ could be continued beyond $T$ from its point at time
$T - \varepsilon/2$, a contradiction. Geodesics are defined for all time, and $g_{\mathbb{D}}$ is complete by Theorem~\ref{riemannian-geometry:theorem-hopf-rinow}.
\end{proof}

\begin{theorem}\label{complex-manifolds:theorem-riemann-surfaces-classification}
Every Riemann surface $X$ is biholomorphic to $\widetilde{X}/\Gamma$, where $\widetilde{X} \in \{\CC\PP^1, \CC, \mathbb{D}\}$ is its universal covering and $\Gamma \cong \pi_1(X)$ is
a group of automorphisms of $\widetilde{X}$ acting freely and properly. Moreover:
\begin{enumerate}
\item $\widetilde{X} = \CC\PP^1$ if and only if $X \cong \CC\PP^1$;
\item $\widetilde{X} = \CC$ if and only if $X$ is biholomorphic to $\CC$, to $\CC^\times$ or to a complex torus $\CC/\Lambda$;
\item in all other cases $\widetilde{X} = \mathbb{D}$, and $X$ carries a unique complete metric of constant curvature $-1$ for which the projection $\mathbb{D} \to X$ is a local
isometry; in holomorphic coordinates it has the form $\lambda(z)|dz|^2$.
\end{enumerate}
\end{theorem}

\begin{proof}
The first statement is Theorem~\ref{complex-manifolds:theorem-uniformisation} and the remark after Lemma~\ref{complex-manifolds:lemma-coverings}. (1) By
Lemma~\ref{complex-manifolds:lemma-automorphisms}(1) every automorphism of $\CC\PP^1$ has a fixed point, so a group acting freely on it is trivial. (2) By
Lemma~\ref{complex-manifolds:lemma-automorphisms}(2), $\Gamma$ consists of translations $z \mapsto z + \lambda$, and $\Lambda = \{\lambda\}$ is a subgroup of $\CC = \RR^2$. It is
discrete: it is the orbit of $0$, which is a fibre of the covering $\CC \to X$. By Lemma~\ref{lie-groups:lemma-discrete-subgroups}, $\Lambda$ is $0$, $\ZZ\omega$ with
$\omega \neq 0$, or a lattice. In the first case $X \cong \CC$; in the second, $z \mapsto e^{2\pi iz/\omega}$ is a covering $\CC \to \CC^\times$ whose fibres are the cosets of
$\ZZ\omega$, so it induces a bijective local biholomorphism $\CC/\ZZ\omega \to \CC^\times$; in the third $X$ is a complex torus. Conversely these surfaces have universal covering
$\CC$. (3) The group $\Gamma \subseteq \mathrm{Aut}(\mathbb{D})$ acts by isometries of $g_{\mathbb{D}}$ (Proposition~\ref{complex-manifolds:proposition-poincare-metric}), so
$g_{\mathbb{D}}$ descends to a metric on $X$ making the projection a local isometry (Proposition~\ref{riemannian-geometry:proposition-quotient-metric}), of curvature $-1$; it is
complete because geodesics of $X$ lift to geodesics of the complete $\mathbb{D}$ (as in the proof of Theorem~\ref{riemannian-geometry:theorem-bonnet-myers}), and in the
holomorphic coordinates given by local inverses of the projection it is $4|dz|^2/(1 - |z|^2)^2$. It is unique because the projection is a local diffeomorphism, so the metric on
$X$ is determined by its pullback $g_{\mathbb{D}}$.
\end{proof}

\begin{corollary}\label{complex-manifolds:corollary-compact-uniformisation}
Let $X$ be a compact Riemann surface of genus $g$ (Theorem~\ref{topological-manifolds:theorem-classification-surfaces}; $X$ is orientable). If $g = 0$ then $X \cong \CC\PP^1$; if
$g = 1$ then $X$ is a complex torus $\CC/\Lambda$; and if $g \geq 2$ then $X \cong \mathbb{D}/\Gamma$, and $X$ carries a metric of constant curvature $-1$, compatible with its
complex structure, with area $4\pi(g - 1)$.
\end{corollary}

\begin{proof}
If $g = 0$, $X$ is simply connected (Corollary~\ref{homotopy:corollary-spheres}), so $X$ is its own universal covering, compact, hence $\CC\PP^1$. If $g \geq 1$, $\pi_1(X)$ is
nontrivial, as its abelianisation $H_1(X; \ZZ)$ (Theorem~\ref{singular-homology:theorem-hurewicz}) is $\ZZ^{2g}$
(Example~\ref{singular-homology:example-cellular-computations}), so $\widetilde{X} \neq \CC\PP^1$ by
Theorem~\ref{complex-manifolds:theorem-riemann-surfaces-classification}(1). If $\widetilde{X} = \mathbb{D}$, the metric of curvature $-1$ has area $A$ with $-A = 2\pi(2 - 2g)$ by
the Gauss--Bonnet theorem (Theorem~\ref{riemannian-geometry:theorem-gauss-bonnet}), so $A = 4\pi(g - 1)$ and $g \geq 2$. If $\widetilde{X} = \CC$, then $X$ is compact among
$\CC$, $\CC^\times$ and the tori, so it is a torus, of genus $1$. Hence $g = 1$ forces $\widetilde{X} = \CC$ and $g \geq 2$ forces $\widetilde{X} = \mathbb{D}$.
\end{proof}

\begin{corollary}[Little Picard theorem]\label{complex-manifolds:corollary-picard}
A nonconstant entire function takes every complex value with at most one exception.
\end{corollary}

\begin{proof}
Suppose the entire function $f$ omits two values; composing with an affine map we may assume they are $0$ and $1$, so $f \colon \CC \to X = \CC \setminus \{0, 1\}$. The fundamental
group of $X$ is not abelian: $X = U \cup V$ with $U = \{\operatorname{Re}z < \frac34\} \setminus \{0\}$ and $V = \{\operatorname{Re}z > \frac14\} \setminus \{1\}$, where
$U \cap V$ is a convex strip, hence simply connected, and $U$, $V$ deformation retract onto the circles of radius $\frac14$ around $0$ and $1$ (radially, along segments
that stay in the convex half planes), so $\pi_1(U) \cong \pi_1(V) \cong \ZZ$ (Lemma~\ref{homotopy:lemma-deformation-retract} and Theorem~\ref{homotopy:theorem-circle}) and
$\pi_1(X) \cong \ZZ * \ZZ$ by Corollary~\ref{homotopy:corollary-spheres}, a non-abelian group (Example~\ref{homotopy:example-wedge}). Since $X$ is not compact and $\pi_1(X)$ is not abelian, while
$\CC$, $\CC^\times$ and the tori have abelian fundamental groups $0$, $\ZZ$ and $\ZZ^2$, Theorem~\ref{complex-manifolds:theorem-riemann-surfaces-classification} gives a covering
map $p \colon \mathbb{D} \to X$ that is a local biholomorphism. As $\CC$ is simply connected, $f$ lifts to a continuous $\tilde f \colon \CC \to \mathbb{D}$ with $p \circ \tilde f = f$
(Proposition~\ref{homotopy:proposition-lifting-criterion}), which is holomorphic, being locally a local inverse of $p$ composed with $f$. Then $\tilde f$ is a bounded entire function,
hence constant (Corollary~\ref{complex-analysis:corollary-liouville}), and so is $f$.
\end{proof}

\begin{remark}\label{complex-manifolds:remark-uniformisation-metrics}
For compact surfaces the three cases match the sign of the Euler characteristic: $\CC\PP^1$ carries the round metric of curvature $1$ (Example~\ref{riemannian-geometry:example-conformal-metrics}(3)),
complex tori the flat metric $|dz|^2$, and surfaces of genus $g \geq 2$ the hyperbolic metric, all compatible with the complex structure; the sign of the curvature is forced by the
Gauss--Bonnet theorem. Schwarz--Pick also shows that the hyperbolic metric is intrinsic: a holomorphic map $f \colon X \to Y$ between Riemann surfaces covered by the disc
lifts to a holomorphic map $\mathbb{D} \to \mathbb{D}$ of the universal coverings (Proposition~\ref{homotopy:proposition-lifting-criterion}), so it does not increase the lengths of
curves for the hyperbolic metrics of Theorem~\ref{complex-manifolds:theorem-riemann-surfaces-classification}(3).
\end{remark}

\section{Almost complex structures and forms of type $(p,q)$}\label{complex-manifolds:section-pq}

\begin{definition}\label{complex-manifolds:definition-almost-complex}
An \emph{almost complex structure} on a smooth manifold $M$ is an endomorphism $J$ of $TM$ with $J^2 = -\id$. It makes each $T_pM$ a complex vector space, and gives the decomposition
$T_pM \otimes \CC = T^{1,0}_p \oplus T^{0,1}_p$ into the $\pm i$ eigenspaces, with $T^{0,1} = \overline{T^{1,0}}$, and dually
$\Lambda^k T^*_pM \otimes \CC = \bigoplus_{p+q=k}\Lambda^{p,q}$ (Proposition~\ref{multilinear-algebra:proposition-pq-decomposition}). A complex manifold has a canonical almost complex structure, given in
holomorphic coordinates $z^j = x^j + iy^j$ by $J\partial_{x^j} = \partial_{y^j}$, $J\partial_{y^j} = -\partial_{x^j}$; then $T^{1,0}$ has frame $\partial/\partial z^j$ and $\Lambda^{p,q}$ has frame
$dz^I \wedge d\bar z^J$ with $|I| = p$, $|J| = q$. An almost complex structure is \emph{integrable} if it arises from a complex structure.
\end{definition}

\begin{definition}\label{complex-manifolds:definition-dbar}
On a complex manifold, $\Omega^{p,q}(X)$ denotes the smooth $(p,q)$-forms, so $\Omega^k(X, \CC) = \bigoplus_{p+q=k}\Omega^{p,q}(X)$. Writing $d = \partial + \bar\partial$ with
$\partial \colon \Omega^{p,q} \to \Omega^{p+1,q}$ and $\bar\partial \colon \Omega^{p,q} \to \Omega^{p,q+1}$ the components of $d$, one has in coordinates
\[
\partial\left(\sum f_{IJ}\,dz^I \wedge d\bar z^J\right) = \sum\frac{\partial f_{IJ}}{\partial z^k}\,dz^k \wedge dz^I \wedge d\bar z^J, \qquad
\bar\partial\left(\sum f_{IJ}\,dz^I \wedge d\bar z^J\right) = \sum\frac{\partial f_{IJ}}{\partial\bar z^k}\,d\bar z^k \wedge dz^I \wedge d\bar z^J .
\]
\end{definition}

\begin{proposition}\label{complex-manifolds:proposition-dbar}
On a complex manifold $d = \partial + \bar\partial$, and $\partial^2 = \bar\partial^2 = \partial\bar\partial + \bar\partial\partial = 0$. A function is holomorphic if and only if $\bar\partial f = 0$, and a $(p, 0)$-form
$\alpha$ is \emph{holomorphic} (has holomorphic coefficients) if and only if $\bar\partial\alpha = 0$; the sheaf of holomorphic $p$-forms is denoted $\Omega^p_X$, with $\Omega^0_X = \mathcal{O}_X$ and
$\Omega^n_X = K_X$ the \emph{canonical bundle}.
\end{proposition}

\begin{proof}
In holomorphic coordinates $d$ of a $(p,q)$-form has components of types $(p+1, q)$ and $(p, q+1)$ only, by the coordinate formula for $d$; these are $\partial$ and $\bar\partial$. Then $0 = d^2$ decomposes into
its components of types $(p+2, q)$, $(p+1, q+1)$ and $(p, q+2)$, which are $\partial^2$, $\partial\bar\partial + \bar\partial\partial$ and $\bar\partial^2$. The statement about functions is the Cauchy--Riemann
equations (Definition~\ref{several-complex-variables:definition-holomorphic}), and for forms it is the same condition on the coefficients.
\end{proof}

\begin{theorem}[Newlander--Nirenberg]\label{complex-manifolds:theorem-newlander-nirenberg}
An almost complex structure $J$ on a smooth manifold is integrable if and only if the \emph{Nijenhuis tensor} vanishes,
\[
N_J(X, Y) = [JX, JY] - J[JX, Y] - J[X, JY] - [X, Y] = 0,
\]
equivalently if and only if $T^{0,1}$ is closed under the Lie bracket, equivalently if and only if $d = \partial + \bar\partial$ on $1$-forms with respect to the induced type decomposition.
\end{theorem}

The implication ``integrable $\Rightarrow$ $N_J = 0$'' is a computation in holomorphic coordinates, where the brackets of the coordinate fields vanish; the converse is a hard analytic theorem, proved in
\cite[Chapter 2]{huybrechts-complex} and, in the real-analytic case, by the Frobenius theorem (Theorem~\ref{smooth-manifolds:theorem-frobenius}) applied to $T^{0,1}$ after complexification.

\begin{proposition}\label{complex-manifolds:proposition-surfaces-almost-complex}
Every almost complex structure on a manifold of real dimension $2$ has vanishing Nijenhuis tensor. Consequently, by Theorem~\ref{complex-manifolds:theorem-newlander-nirenberg}, every oriented surface with a Riemannian metric is a Riemann
surface, with $J$ the rotation by $\pi/2$ in each tangent plane.
\end{proposition}

\begin{proof}
$N_J$ is antisymmetric and $C^\infty$-bilinear, so it suffices to evaluate it on $X$ and $JX$ for a local nowhere vanishing field $X$, which together with $JX$ form a frame. Using $J^2 = -1$,
$N_J(X, JX) = [JX, -X] - J[JX, JX] - J[X, -X] - [X, JX] = [X, JX] - [X, JX] = 0$. On an oriented surface with a metric, rotation by $\pi/2$ in the positive sense is a smooth endomorphism with $J^2 = -1$.
\end{proof}

\section{Analytic subvarieties and blow-ups}\label{complex-manifolds:section-subvarieties}

\begin{definition}\label{complex-manifolds:definition-analytic-subvariety}
An \emph{analytic subvariety} of a complex manifold $X$ is a closed subset that is locally the common zero set of finitely many holomorphic functions (Definition~\ref{several-complex-variables:definition-analytic-set});
it is \emph{smooth}, that is, a complex submanifold, at its regular points. By Theorem~\ref{several-complex-variables:theorem-analytic-structure} the singular locus is a proper analytic subvariety, the regular
locus is dense, and $Z$ has a well-defined dimension on each irreducible component.
\end{definition}

\begin{proposition}\label{complex-manifolds:proposition-subvariety-class}
Let $Z \subseteq X$ be an analytic subvariety of pure complex dimension $k$ in a compact complex manifold of dimension $n$. Then integration over the regular locus,
$\varphi \mapsto \int_{Z_{\mathrm{reg}}}\varphi$, converges for every smooth $2k$-form $\varphi$ and defines a closed current of degree $2n - 2k$; its class in $H^{2n-2k}(X; \RR)$ is the \emph{cycle class} $[Z]$.
Moreover $\int_Z\varphi \geq 0$ when $\varphi$ is a positive $(k,k)$-form.
\end{proposition}

\begin{proof}[Sketch of proof]
Convergence follows from the local structure of analytic sets: near a regular point $Z$ is a submanifold, and near a singular point it is a finite branched cover of a polydisc
(Theorem~\ref{several-complex-variables:theorem-analytic-structure}), so the volume of $Z \cap K$ is finite for compact $K$ (in a chart, Wirtinger's inequality bounds the volume by the integral of a power of the
standard form $\frac{i}{2}\sum_jdz^j \wedge d\bar z^j$ of $\CC^n$). Closedness, that is, Stokes' theorem across the singular locus, is proved by cutting out a neighbourhood of the singular set, whose boundary
has small area because the singular locus has real codimension at least $2$. The positivity statement is local linear algebra. Details are in \cite[Chapter 0, Section 2]{griffiths-harris} and
\cite[Chapter 11]{voisin-hodge-1}.
\end{proof}

\begin{remark}\label{complex-manifolds:remark-cycle-class-nonzero}
If $X$ carries a K\"ahler form $\omega$ (Chapter~\ref{kahler}) and $Z \neq \emptyset$, then $[Z] \neq 0$: pairing with the closed form $\omega^k$ gives $\int_Z\omega^k > 0$, the volume of $Z$ up to a constant.
On non-K\"ahler manifolds cycle classes can vanish; for instance on the Hopf surface of Exercise~\ref{complex-manifolds:exercise-hopf-surface} the elliptic curves $\{z_1 = 0\}/\langle 2 \rangle$ have
zero class, since $H^2$ of the Hopf surface vanishes.
\end{remark}

\begin{definition}\label{complex-manifolds:definition-blowup}
The \emph{blow-up} of $\CC^n$ at the origin is
\[
\mathrm{Bl}_0\CC^n = \{(z, \ell) \in \CC^n \times \CC\PP^{n-1} : z \in \ell\},
\]
with the projection $\pi$ to $\CC^n$. It is a complex manifold, $\pi$ is a biholomorphism over $\CC^n \setminus \{0\}$, and $E = \pi^{-1}(0) \cong \CC\PP^{n-1}$ is the \emph{exceptional divisor}. For a point $p$ of
a complex manifold $X$ one defines $\mathrm{Bl}_pX$ by gluing $\mathrm{Bl}_0\CC^n$ in a chart, and similarly along a complex submanifold $Z \subseteq X$, with $E = \PP(N_{Z/X})$.
\end{definition}

\begin{proposition}\label{complex-manifolds:proposition-blowup}
$\mathrm{Bl}_pX$ is a compact complex manifold if $X$ is, $\pi$ is proper and holomorphic, and
\[
H^k(\mathrm{Bl}_pX; \ZZ) \cong H^k(X; \ZZ) \oplus \bigoplus_{0 < j < n} H^{k-2j}(\mathrm{pt}; \ZZ),
\]
so blowing up a point adds one generator in each even degree $2, 4, \ldots, 2n-2$. The exceptional divisor has normal bundle $\mathcal{O}(-1)$ in $\mathrm{Bl}_0\CC^n$.
\end{proposition}

\begin{proof}
In the chart $\ell = [1 : w_2 : \cdots : w_n]$ the map $(u, w_2, \ldots, w_n) \mapsto (u, uw_2, \ldots, uw_n)$ identifies $\CC^n$ with an open subset of $\mathrm{Bl}_0\CC^n$, which shows that the blow-up is a complex
manifold and that $\pi$ is holomorphic and biholomorphic away from $E$; properness holds because $\CC\PP^{n-1}$ is compact. In these coordinates $E$ is $\{u = 0\}$, and the transition functions of its
normal bundle are those of the tautological bundle, so $N_{E} \cong \mathcal{O}(-1)$. For the cohomology, use the Mayer--Vietoris sequence for $\mathrm{Bl}_pX = (\mathrm{Bl}_pX \setminus E) \cup \pi^{-1}(B)$ with
$B$ a small ball, together with $\pi^{-1}(B) \simeq \CC\PP^{n-1}$ and $\mathrm{Bl}_pX \setminus E \cong X \setminus \{p\}$, and Example~\ref{singular-cohomology:example-cohomology-groups}.
\end{proof}

\begin{remark}\label{complex-manifolds:remark-blowup-birational}
Blow-ups are the basic birational modifications: Hironaka's resolution of singularities states that every analytic (or algebraic, in characteristic zero) variety becomes smooth after finitely many blow-ups
along smooth centres, and the weak factorisation theorem says that any two birational smooth projective varieties are connected by blow-ups and blow-downs. Both are used to reduce statements about the Hodge
conjecture to the smooth projective case (Chapter~\ref{mixed-hodge-structures}).
\end{remark}

\section{Exercises}\label{complex-manifolds:section-exercises}

\begin{exercise}\label{complex-manifolds:exercise-tori}
Show that every complex torus of dimension $1$ is biholomorphic to $\CC/(\ZZ + \tau\ZZ)$ with $\operatorname{Im}\tau > 0$, and that two such are biholomorphic if and only if the parameters are related by
an element of $\SL_2(\ZZ)$ acting by M\"obius transformations.
\end{exercise}

\begin{exercise}\label{complex-manifolds:exercise-projective-space-charts}
Verify that the transition maps of $\CC\PP^n$ are holomorphic and that $\CC\PP^1$ is biholomorphic to the Riemann sphere $\CC \cup \{\infty\}$, with $S^2$ as underlying smooth manifold.
\end{exercise}

\begin{exercise}\label{complex-manifolds:exercise-hopf-surface}
The \emph{Hopf surface} is $X = (\CC^2 \setminus \{0\})/\langle z \mapsto 2z\rangle$. Show that it is a compact complex manifold diffeomorphic to $S^1 \times S^3$, and deduce that $b_1(X) = 1$; conclude that $X$
carries no K\"ahler metric (Chapter~\ref{hodge-decomposition}), since odd Betti numbers of compact K\"ahler manifolds are even.
\end{exercise}

\begin{solution}
The action is free and properly discontinuous, so the quotient is a complex manifold, and it is compact because the annulus $1 \leq |z| \leq 2$ maps onto it. The diffeomorphism with $S^1 \times S^3$ comes from
the polar decomposition $\CC^2 \setminus \{0\} \cong S^3 \times \RR_{>0}$, with the action becoming translation on the second factor. Then $b_1(S^1 \times S^3) = 1$ by the K\"unneth formula.
\end{solution}

\begin{exercise}\label{complex-manifolds:exercise-almost-complex-s6}
Show that $S^2$ and $S^6$ admit almost complex structures (for $S^6$ use the octonionic cross product), and that $S^2$'s is integrable.
\end{exercise}

\begin{exercise}\label{complex-manifolds:exercise-residue-infinity}
Compute $\int_{|z| = 2}\frac{z^5}{z^6 - 1}\,dz$ and $\int_{|z| = 2}\frac{dz}{(z - 1)^3(z^2 + 9)}$ (the circle traversed once counterclockwise) by computing a single residue at $\infty$ on $\CC\PP^1$ instead of the residues
inside the circle.
\end{exercise}

\begin{solution}
By Theorem~\ref{complex-manifolds:theorem-global-residue} on $\CC\PP^1$, the integral over the circle, which is $2\pi i$ times the sum of the residues inside it, equals $-2\pi i$ times the sum of the residues
outside it, including $\infty$. For $\omega = z^5\,dz/(z^6 - 1)$ all six poles lie inside. With $w = 1/z$, $\frac{z^5}{z^6 - 1} = \frac{w}{1 - w^6}$ and $dz = -dw/w^2$, so $\omega = -\frac{dw}{w(1 - w^6)}$ has residue $-1$
at $\infty$, and the integral is $2\pi i$. (Directly: each pole $\zeta$ is simple with residue $\zeta^5/(6\zeta^5) = \frac16$.) For the second form the poles $\pm3i$ lie outside the circle and $1$ inside. At
$\infty$ the form is $O(w^3)\,dw$, holomorphic, so the integral is $-2\pi i(\operatorname{Res}_{3i} + \operatorname{Res}_{-3i})$, with
$\operatorname{Res}_{\pm3i} = \frac{1}{(\pm3i - 1)^3(\pm6i)}$. Now $(3i - 1)^3 = 26 - 18i$ and $(3i - 1)^3 \cdot 6i = 108 + 156i$, so $\operatorname{Res}_{3i} = \frac{1}{108 + 156i}$ and
$\operatorname{Res}_{-3i} = \frac{1}{108 - 156i}$ is its conjugate. Their sum is $\frac{216}{108^2 + 156^2} = \frac{216}{36000} = \frac{3}{500}$, and the integral is $-\frac{3\pi i}{250}$. (Directly: the
residue at the triple pole $1$ is $\frac12\frac{d^2}{dz^2}(z^2 + 9)^{-1}$ at $z = 1$, which is $\frac12\bigl(-\frac{2}{100} + \frac{8}{1000}\bigr) = -\frac{3}{500}$.)
\end{solution}

\begin{exercise}\label{complex-manifolds:exercise-abel-torus}
Let $f$ be a nonconstant meromorphic function on $X = \CC/\Lambda$, $\Lambda = \ZZ\omega_1 + \ZZ\omega_2$, viewed as a $\Lambda$-periodic meromorphic function on $\CC$. Let $a_1, \ldots, a_k$ and $b_1, \ldots, b_k$ be
its zeros and poles in a fundamental parallelogram $P$ with vertices $a$, $a + \omega_1$, $a + \omega_1 + \omega_2$, $a + \omega_2$, listed with multiplicity, with $a$ chosen so that no zero or pole lies on
$\partial P$. Show that $\sum_ia_i - \sum_ib_i \in \Lambda$, by computing $\frac{1}{2\pi i}\int_{\partial P}z\frac{f^{\prime}(z)}{f(z)}\,dz$ in two ways. (This is one direction of Abel's theorem for complex tori.)
\end{exercise}

\begin{solution}
At a zero or pole $c$ of order $m$ (negative for a pole), $f^{\prime}/f = m/(z - c) + (\text{holomorphic})$, so $zf^{\prime}/f$ has residue $mc$. The residue theorem in the plane
(Theorem~\ref{complex-analysis:theorem-residues}) gives $\frac{1}{2\pi i}\int_{\partial P}zf^{\prime}/f\,dz = \sum_ia_i - \sum_ib_i$. On the other hand, pair opposite sides: the side from $a + \omega_2$ to
$a + \omega_1 + \omega_2$ is traversed backwards, and on it $f^{\prime}/f$ takes the same values as on the side from $a$ to $a + \omega_1$, so the two sides contribute
\[
\int_a^{a+\omega_1}\bigl(z - (z + \omega_2)\bigr)\frac{f^{\prime}(z)}{f(z)}\,dz = -\omega_2\int_a^{a+\omega_1}\frac{f^{\prime}(z)}{f(z)}\,dz = -2\pi i\,n_1\omega_2,
\]
where $\frac{1}{2\pi i}\int_a^{a+\omega_1}f^{\prime}/f\,dz = n_1$ is the winding number around $0$ of the closed curve $t \mapsto f(a + t\omega_1)$, $t \in [0, 1]$ (closed because $f$ is periodic), an
integer. In the same way the other two sides contribute $2\pi i\,n_2\omega_1$ with $n_2 \in \ZZ$. So $\sum_ia_i - \sum_ib_i = n_2\omega_1 - n_1\omega_2 \in \Lambda$. Applying the same computation to
$f^{\prime}/f$ instead of $zf^{\prime}/f$ gives $k$ zeros and $k$ poles again, in accordance with Corollary~\ref{complex-manifolds:corollary-zeros-poles}.
\end{solution}
